


\section{Human Annotation Study}\label{sec:human-annotation}

% As discussed in \S~\ref{sec:dataset/overview}, considering only the cited candidates as the gold-standard can lead to inconsistencies due to variations in how different lawyers perceive cases. To ground our research in reality, we took help from 6 senior law students (pursuing LL.M.) from The West Bengal National University of Juridical Sciences (a very reputed Law school of India) to conduct annotation exercises, under the guidance of a senior faculty member from the same school. The annotators were asked to perform the tasks of PCR and LSR manually. 

% %Additionally, annotation and evaluation by human experts is considered an important step to ground legal NLP research~\citep{chalkidis2021echr,paul2024legal}.Hence, we employed 6 senior law students (pursuing LL.M.) from The West Bengal National University of Juridical Sciences (a very reputed Law school of India) to conduct annotation exercises, under the guidance of a senior faculty member from the same school.

% \noindent\underline{\textbf{Annotation Setup:}} We randomly selected a set of 60 documents (queries) from the test set ($\sim$10\% of total test set). Each student was provided 30 queries, and they were tasked with finding out \emph{all} relevant candidates (both statutes and precedents) from the same retrieval pools.  We created a specific search tool for each retrieval task which, for a given query, returned a ranked list of all candidates (only those belonging to our candidate sets) containing any keyword in the query. This tool works like a simple search engine over only the candidate pools used in our work.
% The annotators were asked to try out an extensive set of keywords for each query (as many as they thought necessary), and mark out the candidates they thought were relevant to the query. After this exercise, each query received three distinct sets of annotated precedents/statutes from three annotators.

% \noindent\underline{\textbf{Salient Statistics:}} For these 60 queries, the average number of cited statutes and precedents is 4.57 and 2.72 respectively, while the average number of human-annotated statutes and precedents are 5.0 and 5.87 respectively. This demonstrates the subjectivity in relevance of candidates for the PCR task in particular.

% \noindent\underline{\textbf{Experiments:}} For these 60 documents, we take some of the best performing methods, and calculate the F1 scores based on -— (i)~only the cited candidates as gold-standard (same setup as Section~\ref{sec:expt}), and (ii)~considering the union of cited candidates and the three sets of annotated precedents/statutes as gold-standard. The results are shown in Table~\ref{tab:human-results}. We observe that the scores are higher with the gold-standard of citations + annotations compared to citations only, for all approaches. This shows that many top-ranked candidates (for every model) are considered relevant by the annotators, even if they are not cited. 
% Importantly, the trends for the models' performances remain broadly the same across both kinds of gold-standard, with the second-stage LLM re-ranking performing the best for both statutes and precedents. This suggests that cited candidates can serve as a good proxy for the set of all possible relevant candidates.

As discussed in \S~\ref{sec:dataset/overview}, using only cited candidates as gold-standard can lead to inconsistencies, since relevance judgments may vary across lawyers based on their individual perspectives. To ground our evaluation in practice, we conducted an annotation study with six senior L.L.M. students from the WB National University of Juridical Sciences (a reputed Indian Law school), under the supervision of a senior faculty member (more details regarding the annotators in App.~\ref{app:annot}). The annotators were asked to manually perform both LSR and PCR tasks.  

\vspace{1mm}
\noindent\underline{\textbf{Annotation Setup:}} We randomly selected 60 queries (about 10\% of the test set). Each query was assigned to three annotators, with each annotator handling 30 queries. To support retrieval, we built a simple search tool over the candidate pools for statutes and precedents, returning ranked lists of candidates containing any keyword from the query. Annotators were asked to try multiple keywords for each query, and mark all candidates they judged relevant. Thus, each query received three independent sets of annotations for both tasks.  

\vspace{1mm}
\noindent\underline{\textbf{Salient Statistics:}} On average, the cited set contained 4.57 statutes and 2.72 precedents per query. In contrast, annotators identified 5.0 statutes and 5.87 precedents. This indicates that statute relevance is relatively stable, with only a small increase over citations, whereas precedent relevance is highly subjective, with many more cases judged relevant than those actually cited.  

\vspace{1mm}
\noindent\underline{\textbf{Results:}} We compared the performance of selected models under two gold standards: (i)~citations only and (ii)~the union of citations and annotations (Table~\ref{tab:human-results}). As expected, scores increase across all models under the broader gold standard, since many top-ranked candidates judged relevant by models were also marked relevant by annotators even if not cited. Importantly, the relative performance trends remain consistent across both setups. For statutes, the stability of annotations ensures that \textit{Stage-2 LLM re-ranking consistently outperforms all other methods}. For precedents, the gains over ensembles are smaller under the annotated gold, likely because lexical-heavy ensembles benefit more directly from the broader definition of relevance. Nevertheless, Stage-2 prompting edges ahead, showing that \textit{LLMs not only outperform ensembles but also preserve their advantage under more realistic evaluation}. Thus, while citations capture only a part of the relevant set, they serve as a sufficiently reliable proxy for system evaluation.


